\documentclass[12pt,a4paper]{article}

% ══════════════════════════════════════════════════════════════════════════════
%  INFORMATIONS À RENSEIGNER AVANT ENVOI
%  Les champs marqués \acompleter{...} apparaissent en rouge dans le PDF :
%  ils signalent ce que seul l'auteur ou le tuteur peut renseigner.
% ══════════════════════════════════════════════════════════════════════════════
\newcommand{\nouncleve}{HOUNKPETOHOU Marius}
\newcommand{\prenomeleve}{Marius}
\newcommand{\numetudiant}{à compléter}
\newcommand{\anneeacad}{2025 -- 2026}
\newcommand{\entreprise}{Choubel Consulting}
\newcommand{\villestage}{Casablanca, Maroc}
\newcommand{\datedebut}{13 juillet 2026}
\newcommand{\datefin}{21 août 2026}
\newcommand{\tuteurnom}{à compléter}
\newcommand{\tuteurfonction}{à compléter}
\newcommand{\binome}{DANSOU René}

% ══════════════════════════════════════════════════════════════════════════════
% Préambule du rapport de stage académique (École Centrale Casablanca)
%
% Deux exigences se superposent ici, et il a fallu les tenir ensemble :
%
%   1. le guide de rédaction de l'école — police 12, interligne 1,5, texte
%      justifié, pagination ;
%   2. la charte des livrables du projet — c'est le gabarit des rapports
%      hebdomadaires remis au cabinet (rapports/_preambule.tex) qui est repris
%      ici : même page de garde sur fond navy, mêmes bandeaux de partie à
%      pastille numérotée, même table des matières, mêmes en-têtes.
%
% Ce qui a été adapté : la typographie du corps (12 pt et interligne 1,5 au lieu
% de 11 pt serré), les marges, la page de garde qui doit porter les mentions
% imposées (numéro d'étudiant, tuteur, dates), et un bandeau sans numéro pour
% les sections non numérotées (remerciements, introduction, conclusion…).
% ══════════════════════════════════════════════════════════════════════════════

% ─── Encodage & langue ────────────────────────────────────────────────────────
\usepackage[T1]{fontenc}
\usepackage[utf8]{inputenc}
\usepackage[french]{babel}
% Pas de lmodern : les rapports hebdomadaires composent en Computer Modern (les
% fontes T1 « cm-super »), et le rapport doit être dans la même police qu'eux.

% ─── Mise en page imposée par le guide ────────────────────────────────────────
\usepackage[left=2.5cm, right=2.5cm, top=2.4cm, bottom=2.5cm]{geometry}
\usepackage{setspace}
\onehalfspacing                     % interligne 1,5
\usepackage{ragged2e}
\justifying                         % texte justifié
\setlength{\parindent}{1.5em}    % même retrait d'alinéa que les rapports hebdo
\setlength{\parskip}{0.22em plus 0.08em minus 0.04em}
\usepackage{indentfirst}
\setlength{\headheight}{15pt}
\setlength{\headsep}{18pt}

% ─── Mathématiques, graphiques & tableaux ─────────────────────────────────────
\usepackage{amsmath,amssymb}
\usepackage{graphicx}
\graphicspath{{img/}{img/photos/}{../}}
\usepackage{float}
\usepackage{booktabs}
\usepackage{array}
\usepackage{tabularx}
\usepackage{multirow}
\usepackage{xcolor}
\usepackage{colortbl}
\usepackage[font=small,labelfont={bf,color=eccblue}]{caption}
\usepackage{tikz}
\usetikzlibrary{arrows.meta,decorations.pathmorphing,patterns,calc,fadings,
	positioning,shapes.geometric}
\usepackage[most]{tcolorbox}
\usepackage{enumitem}
\setlist[itemize,1]{label=\textbullet, leftmargin=1.6em, itemsep=0.22em, topsep=0.32em}
\setlist[enumerate,1]{leftmargin=1.8em, itemsep=0.22em, topsep=0.32em}

% ─── Couleurs (charte du projet, identique aux rapports hebdomadaires) ────────
\definecolor{eccblue}{RGB}{0,82,145}
\definecolor{linkblue}{RGB}{0,0,250}
\definecolor{navy}{RGB}{12,25,48}
\definecolor{accent}{RGB}{0,188,170}
\definecolor{smoke}{RGB}{160,170,185}
\definecolor{darkgray}{RGB}{80,80,80}
\definecolor{choubelor}{RGB}{249,204,40}
\definecolor{arelire}{RGB}{176,0,32}

\newcommand{\coverSerif}{\fontfamily{ptm}\selectfont}
\newcommand{\coverSans}{\fontfamily{phv}\selectfont}

% ─── Liens ────────────────────────────────────────────────────────────────────
\usepackage{hyperref}
\usepackage{bookmark}
\hypersetup{colorlinks=true, linkcolor=linkblue, urlcolor=linkblue,
	citecolor=linkblue, filecolor=linkblue}

% ═══ Bandeaux de partie ═══════════════════════════════════════════════════════
% Repris des rapports hebdomadaires : pastille dégradée portant le numéro,
% barre bleue, titre centré en italique gras, filet de séparation. Les
% coordonnées sont exprimées pour une justification de 16 cm (marges 2,5 cm) ;
% le bandeau déborde de 1,1 cm dans la marge de gauche, comme dans le gabarit
% d'origine, pour que la pastille « sorte » du bloc de texte.
\newcounter{manualbookmark}
\makeatletter

\newcommand{\bandeautitre}[1]{%
	\node[text width=13.6cm,align=center,
	font=\bfseries\itshape\fontsize{17.5}{21}\selectfont,color=black]
	at (9.35,-0.92) {#1};
	\draw[eccblue!55,line width=0.50pt] (1.55,-1.80) -- (17.10,-1.80);}

\newcommand{\mainsection}[1]{%
	\clearpage
	\thispagestyle{empty}%
	\refstepcounter{section}%
	\setcounter{subsection}{0}%
	\stepcounter{manualbookmark}%
	\phantomsection%
	\pdfbookmark[1]{\thesection\space #1}{sec-\themanualbookmark}%
	\addcontentsline{toc}{section}{\protect\numberline{\thesection}#1}%
	\markboth{\thesection. #1}{\thesection. #1}%
	\vspace*{0.05cm}%
	\noindent\hspace*{-1.55cm}%
	\begin{tikzpicture}[x=1cm,y=1cm]
		\shade[inner color=eccblue!18,outer color=eccblue!82]
		(0.85,0.05) circle (1.30);
		\node[anchor=west,font=\fontsize{8.5}{10}\selectfont,color=white]
		at (0.16,0.90) {Partie};
		\node[font=\bfseries\fontsize{32}{32}\selectfont,color=white]
		at (0.87,0.02) {\thesection};
		\fill[eccblue!80] (1.62,-0.16) rectangle (17.10,0.06);
		\bandeautitre{#1}
	\end{tikzpicture}\par
	\vspace{0.70cm}%
}

% Même bandeau, sans numéro : remerciements, glossaire, introduction,
% conclusion, sources, annexes. La pastille passe en teinte d'accent et reste
% muette — écrire un numéro là où il n'y en a pas serait un contresens.
\newcommand{\mainsectionstar}[1]{%
	\clearpage
	\thispagestyle{empty}%
	\stepcounter{manualbookmark}%
	\phantomsection%
	\pdfbookmark[1]{#1}{sec-\themanualbookmark}%
	\addcontentsline{toc}{section}{#1}%
	\markboth{#1}{#1}%
	\vspace*{0.05cm}%
	\noindent\hspace*{-1.55cm}%
	\begin{tikzpicture}[x=1cm,y=1cm]
		\shade[inner color=accent!14,outer color=accent!70]
		(0.85,0.05) circle (1.30);
		\fill[white,opacity=0.85] (0.85,0.05) circle (0.20);
		\fill[accent!80] (1.62,-0.16) rectangle (17.10,0.06);
		\bandeautitre{#1}
	\end{tikzpicture}\par
	\vspace{0.70cm}%
}
\makeatother

% ─── Titres de niveau inférieur ───────────────────────────────────────────────
\usepackage{titlesec}
\titleformat{\subsection}
{\color{eccblue}\large\bfseries}{\thesubsection.}{0.55em}{}
\titleformat{\subsubsection}
{\color{navy}\normalsize\bfseries\itshape}{\thesubsubsection.}{0.5em}{}
\titlespacing*{\subsection}{0pt}{1.25em}{0.45em}
\titlespacing*{\subsubsection}{0pt}{0.95em}{0.35em}

% ─── En-têtes et pieds de page (gabarit hebdomadaire) ─────────────────────────
\usepackage{fancyhdr}
\pagestyle{fancy}
\fancyhf{}
\fancyhead[L]{\scriptsize\leftmark}
\fancyhead[R]{\scriptsize\thepage}
\fancyfoot[L]{\fontsize{8}{9}\selectfont\color{darkgray}\textit{\nouncleve}}
\fancyfoot[C]{\fontsize{9}{10}\selectfont\bfseries\color{eccblue}\thepage}
\fancyfoot[R]{\fontsize{8}{9}\selectfont\color{darkgray}\textit{Stage opérateur --- \entreprise}}
\renewcommand{\headrulewidth}{1.0pt}
\renewcommand{\footrulewidth}{0.8pt}

% ─── Table des matières ───────────────────────────────────────────────────────
\usepackage{tocloft}
\renewcommand{\contentsname}{Table des matières}
\renewcommand{\cfttoctitlefont}{\hfill\LARGE\bfseries\itshape\color{black}}
\renewcommand{\cftaftertoctitle}{\hfill\par\vspace{0.85cm}}
\renewcommand{\cftsecfont}{\small\color{eccblue}\bfseries}
\renewcommand{\cftsecpagefont}{\small\bfseries\color{black}}
\renewcommand{\cftsubsecfont}{\small\color{navy}}
\renewcommand{\cftdotsep}{2.5}

\makeatletter
\newcommand{\printstyledtoc}{%
	\clearpage
	\pagestyle{empty}%
	\thispagestyle{empty}%
	\phantomsection%
	\pdfbookmark[1]{Table des matières}{toc}%
	\vspace*{0.58cm}%
	\noindent
	\begin{tikzpicture}[baseline]
		\fill[gray!28] (0,0.09) rectangle (0.62\linewidth,0.50);
		\node[anchor=east,font=\LARGE\bfseries\itshape,color=black]
		at (\linewidth,0.30) {Table des matières};
		\draw[black,line width=0.55pt] (0,0.00) -- (\linewidth,0.00);
	\end{tikzpicture}\par
	\vspace{0.82cm}%
	{\hypersetup{linkcolor=linkblue,linktoc=section}\@starttoc{toc}}%
	\clearpage
	\pagestyle{fancy}}%
\makeatother

% ─── Encadrés ─────────────────────────────────────────────────────────────────
\newtcolorbox{encadre}[1][]{colback=eccblue!5, colframe=eccblue!55,
	boxrule=0.7pt, arc=2pt, left=9pt, right=9pt, top=7pt, bottom=7pt, #1}
\newtcolorbox{encadrecle}[1][]{colback=accent!7, colframe=accent!70,
	boxrule=0.7pt, arc=2pt, left=9pt, right=9pt, top=7pt, bottom=7pt, #1}
\newtcolorbox{citationbox}[1][]{colback=navy, colframe=navy, coltext=white,
	boxrule=0pt, arc=2pt, left=11pt, right=11pt, top=9pt, bottom=9pt, #1}

% ─── Marqueur des informations restant à compléter ────────────────────────────
% Rend visible, à la relecture, tout ce que seul l'auteur ou le tuteur peut
% renseigner. À supprimer avant l'envoi du rapport.
\newcommand{\acompleter}[1]{{\color{arelire}\bfseries[à compléter : #1]}}

% ─── Insertion d'images tolérante ─────────────────────────────────────────────
\newcommand{\safelogo}[2]{%
	\IfFileExists{#1}{\includegraphics[height=#2,keepaspectratio]{#1}}{\phantom{\rule{0pt}{#2}}}}

% ─── Figure standard, avec crédit optionnel ───────────────────────────────────
% #1 largeur · #2 fichier · #3 légende · #4 crédit (peut être vide)
% La hauteur est bornée : une capture d'écran de page longue dépasserait sinon
% la hauteur du bloc de texte, et LaTeX la repousserait en laissant une page à
% moitié vide.
\newcommand{\figurecredit}[4]{%
	\begin{figure}[H]
		\centering
		\includegraphics[width=#1\linewidth,height=0.74\textheight,keepaspectratio]{#2}
		\caption{#3\ifx\relax#4\relax\else\\{\footnotesize\color{darkgray}\textit{#4}}\fi}
	\end{figure}}

% ═══ Page de garde ════════════════════════════════════════════════════════════
% Reprise de \couverture des rapports hebdomadaires : fond navy pleine page,
% dégradé haut, logos aux deux angles, pastille de titre, titre serif blanc,
% sous-titre en teinte d'accent, blocs d'information étiquetés, ondes en pied
% de page et filet de bas de page. S'y ajoutent les mentions imposées par
% l'école : numéro d'étudiant, tuteur, dates et année académique.
% #1 badge · #2 titre · #3 sous-titre
\newcommand{\couverturestage}[3]{%
	\hypersetup{pageanchor=false}%
	\begin{titlepage}
		\thispagestyle{empty}
		\begin{tikzpicture}[remember picture,overlay]
			\fill[navy] (current page.south west) rectangle (current page.north east);
			\fill[navy,opacity=0.35]
			(current page.north west)
			rectangle ([yshift=14.8cm]current page.south west -| current page.east);
			\fill[navy,path fading=north]
			([yshift=14.8cm]current page.south west)
			rectangle ([yshift=21cm]current page.south west -| current page.east);

			% Logos : école à gauche, organisme d'accueil à droite
			\node[anchor=north west, inner sep=0pt, fill=white, rounded corners=3pt,
			inner xsep=7pt, inner ysep=7pt]
			at ([xshift=2.2cm,yshift=-1.9cm]current page.north west)
			{\safelogo{../logo.jpg}{2.2cm}};
			\node[anchor=north east, inner sep=0pt, fill=white, rounded corners=3pt,
			inner xsep=7pt, inner ysep=7pt]
			at ([xshift=-2.2cm,yshift=-1.9cm]current page.north east)
			{\safelogo{../app/static/img/logo-choubel.jpg}{2.2cm}};

			% Année académique
			\node[anchor=north east, font=\coverSans\fontsize{8.5}{11}\selectfont, text=smoke]
			at ([xshift=-2.2cm,yshift=-5.1cm]current page.north east)
			{Année académique \anneeacad};

			% Pastille de titre
			\node[anchor=south west,
			font=\coverSerif\fontsize{10.5}{12}\selectfont\bfseries,
			text=white, fill=eccblue, fill opacity=0.85, text opacity=1,
			inner xsep=6pt, inner ysep=4pt, rounded corners=2pt]
			at ([xshift=2.2cm,yshift=19.9cm]current page.south west)
			{#1};

			% Titre et sous-titre
			\node[anchor=south west, text width=16.6cm,
			font=\coverSerif\bfseries\fontsize{24}{29}\selectfont, text=white, inner sep=0pt]
			at ([xshift=2.2cm,yshift=14.9cm]current page.south west)
			{\raggedright\hyphenpenalty=10000\exhyphenpenalty=10000 #2};
			\node[anchor=south west, text width=16.6cm,
			font=\coverSerif\fontsize{14}{18}\selectfont, text=accent, inner sep=0pt]
			at ([xshift=2.2cm,yshift=13.1cm]current page.south west)
			{\raggedright #3};
			\draw[accent!60, line width=1.1pt]
			([xshift=2.2cm,yshift=12.5cm]current page.south west)
			-- ([xshift=6.2cm,yshift=12.5cm]current page.south west);

			% Colonne de gauche : élève, puis période
			\node[anchor=north west, font=\coverSans\fontsize{7.5}{9.5}\selectfont, text=accent]
			at ([xshift=2.2cm,yshift=11.7cm]current page.south west) {ÉLÈVE INGÉNIEUR};
			\node[anchor=north west, text width=8cm,
			font=\coverSerif\fontsize{12.5}{16}\selectfont, text=white, inner sep=0pt]
			at ([xshift=2.2cm,yshift=11.0cm]current page.south west)
			{\textbf{HOUNKPETOHOU} \prenomeleve \\[2pt]
				{\fontsize{9.5}{12}\selectfont\color{smoke}
					N\textsuperscript{o} d'étudiant : \numetudiant}};

			\node[anchor=north west, font=\coverSans\fontsize{7.5}{9.5}\selectfont, text=accent]
			at ([xshift=2.2cm,yshift=8.5cm]current page.south west) {PÉRIODE DE STAGE};
			\node[anchor=north west, text width=8cm,
			font=\coverSerif\fontsize{11.5}{15}\selectfont, text=white, inner sep=0pt]
			at ([xshift=2.2cm,yshift=7.8cm]current page.south west)
			{Du \datedebut{} au \datefin \\[2pt]
				{\fontsize{9.5}{12}\selectfont\color{smoke} Soit six semaines --- \villestage}};

			% Filet vertical, comme sur la couverture des rapports hebdomadaires
			\draw[smoke!30,line width=0.4pt]
			([xshift=10.4cm,yshift=11.9cm]current page.south west)
			-- ([xshift=10.4cm,yshift=6.9cm]current page.south west);

			% Colonne de droite : organisme d'accueil, puis tuteur
			\node[anchor=north west, font=\coverSans\fontsize{7.5}{9.5}\selectfont, text=accent]
			at ([xshift=11.2cm,yshift=11.7cm]current page.south west) {ORGANISME D'ACCUEIL};
			\node[anchor=north west, text width=6.4cm,
			font=\coverSerif\fontsize{12.5}{16}\selectfont, text=white, inner sep=0pt]
			at ([xshift=11.2cm,yshift=11.0cm]current page.south west)
			{\textbf{\entreprise} \\[2pt]
				{\fontsize{9.5}{12}\selectfont\color{smoke} \villestage}};

			\node[anchor=north west, font=\coverSans\fontsize{7.5}{9.5}\selectfont, text=accent]
			at ([xshift=11.2cm,yshift=8.5cm]current page.south west) {TUTEUR ORGANISME D'ACCUEIL};
			\node[anchor=north west, text width=6.4cm,
			font=\coverSerif\fontsize{11.5}{15}\selectfont, text=white, inner sep=0pt]
			at ([xshift=11.2cm,yshift=7.8cm]current page.south west)
			{\tuteurnom \\[2pt]
				{\fontsize{9.5}{12}\selectfont\color{smoke} \tuteurfonction}};

			% Ondes de pied de page : le motif des couvertures du projet
			\begin{scope}[shift={([xshift=10.5cm,yshift=3.4cm]current page.south west)},
				accent,line width=0.7pt,opacity=0.22]
				\foreach \a/\p in {0.6/0, 0.45/40, 0.75/80, 0.35/120} {
					\draw plot[domain=-4.5:4.5,samples=80,smooth]
					(\x,{\a*sin(200*\x+\p)});
				}
			\end{scope}

			\draw[smoke!25,line width=0.3pt]
			([xshift=2.2cm,yshift=1.8cm]current page.south west)
			-- ([xshift=-2.2cm,yshift=1.8cm]current page.south east);
			\node[anchor=south west, font=\coverSans\fontsize{7.5}{9}\selectfont, text=smoke!60]
			at ([xshift=2.2cm,yshift=1.0cm]current page.south west)
			{École Centrale Casablanca --- Rapport de stage opérateur de première année
				--- Projet mené en binôme avec \binome};
		\end{tikzpicture}
	\end{titlepage}
	\hypersetup{pageanchor=true}%
}


\hypersetup{pdftitle={Rapport de stage opérateur -- \nouncleve},
	pdfauthor={\nouncleve}}

\begin{document}

% ══════════════════════════════════════════════════════════════════════════════
%  PAGE DE GARDE — gabarit des rapports hebdomadaires du projet
% ══════════════════════════════════════════════════════════════════════════════
\couverturestage
	{RAPPORT DE STAGE OPÉRATEUR --- PREMIÈRE ANNÉE}
	{Concevoir un outil d'aide à la décision\\ pour un cabinet de conseil\\ en investissement immobilier}
	{L'application web « RentImmo » --- de la formule financière\\ au geste métier du conseiller}

\newpage
\pagenumbering{roman}
\setcounter{page}{1}

% ══════════════════════════════════════════════════════════════════════════════
\mainsectionstar{Remerciements}

Je tiens d'abord à remercier \tuteurnom, \tuteurfonction{} au sein de
\entreprise, pour m'avoir accueilli au sein du cabinet et pour avoir accepté de
consacrer du temps à répondre, en détail et sans complaisance, aux questions que
nous lui avons adressées sur sa méthode de travail. Ses réponses ont fait plus
que compléter notre cahier des charges : elles en ont corrigé les prémisses. Le
tournant de ce stage tient à cet échange, et je lui suis reconnaissant de la
franchise avec laquelle il a décrit un métier bien plus nuancé que le modèle que
nous en avions formé.

Je remercie également l'ensemble des collaborateurs du cabinet
\acompleter{citer nommément les personnes rencontrées} pour leur disponibilité
et pour les explications concrètes qu'ils ont bien voulu me donner sur le
déroulement réel d'un dossier client.

Je remercie \binome, avec qui j'ai mené ce projet en binôme pendant six
semaines. Travailler à deux sur une même base de code impose une discipline —
écrire avant de coder, tester avant de livrer, expliquer avant de fusionner —
dont j'ai mesuré la valeur au fil des semaines.

Enfin, je remercie l'équipe pédagogique de l'École Centrale Casablanca, et en
particulier \acompleter{nom du responsable des stages}, pour l'encadrement et le
cadre méthodologique de ce stage opérateur.

% ══════════════════════════════════════════════════════════════════════════════
\mainsectionstar{Glossaire et abréviations}

% Un tableau par famille de vocabulaire : une seule tabularx dépasserait la
% hauteur d'une page, et LaTeX ne sait pas la couper.
\newcommand{\glossaire}[2]{%
	\begin{table}[H]
		\centering
		\small
		\setstretch{1.05}
		\renewcommand{\arraystretch}{1.40}
		\begin{tabular}{|>{\bfseries\raggedright\arraybackslash}p{3.7cm}|>{\raggedright\arraybackslash}p{11.1cm}|}
			\hline
			\rowcolor{eccblue!15}
			\multicolumn{2}{|l|}{\textbf{\color{eccblue}#1}} \\
			\hline
			#2
		\end{tabular}
	\end{table}}

\glossaire{Vocabulaire métier}{%
			\rowcolor{white}
			VEFA & Vente en l'état futur d'achèvement, ou « achat sur plan » : le bien
			est acheté avant sa construction et livré plusieurs mois, voire plusieurs
			années, après la signature. \\
			\hline
			\rowcolor{gray!8}
			Compromis & Avant-contrat par lequel vendeur et acheteur s'engagent, avant
			le passage devant notaire. \\
			\hline
			\rowcolor{white}
			Portage foncier & Acquisition d'un terrain conservé sans mise en valeur,
			dans l'attente d'une valorisation liée à l'évolution de son environnement. \\
			\hline
			\rowcolor{gray!8}
			Standing & Niveau de gamme d'un bien : économique, social, moyen standing,
			haut standing, luxe. \\
			\hline
			\rowcolor{white}
			Rendement brut & Loyer annuel rapporté au coût total d'acquisition. \\
			\hline
			\rowcolor{gray!8}
			Rendement net & Rendement après charges, vacance locative et frais de
			gestion, avant impôt. \\
			\hline
			\rowcolor{white}
			Rendement net-net & Rendement net après imposition des revenus locatifs. \\
			\hline
			\rowcolor{gray!8}
			Cash-flow & Trésorerie effectivement encaissée (ou déboursée) sur une
			période, une fois les annuités d'emprunt payées. \\
			\hline
			\rowcolor{white}
			Vacance locative & Part de l'année pendant laquelle le bien n'est pas loué. \\
			\hline
			\rowcolor{gray!8}
			Plus-value & Écart entre le prix de revente d'un bien et son prix d'achat. \\
			\hline
			\rowcolor{white}
			Valeur créée & Ce que l'opération a rapporté au client sur son horizon,
			revente comprise, une fois tout ce qu'il a déboursé retranché. \\
			\hline
}

\glossaire{Vocabulaire financier}{%
			\rowcolor{white}
			VAN & Valeur actuelle nette : somme des flux futurs ramenés à
			aujourd'hui par un taux d'actualisation. \\
			\hline
			\rowcolor{gray!8}
			TRI & Taux de rentabilité interne : taux d'actualisation qui annule la VAN,
			interprété comme le rendement annuel de l'argent investi. \\
			\hline
			\rowcolor{white}
			CRD & Capital restant dû sur un emprunt à une date donnée. \\
			\hline
			\rowcolor{gray!8}
			Horizon & Durée sur laquelle le client projette de conserver le bien. \\
			\hline
			\rowcolor{white}
			MAD & Dirham marocain, devise de référence de l'outil. \\
			\hline
}

\glossaire{Vocabulaire technique}{%
			\rowcolor{white}
			Flask & Cadriciel (\textit{framework}) web écrit en Python, utilisé pour
			développer l'application. \\
			\hline
			\rowcolor{gray!8}
			Blueprint & Module d'une application Flask regroupant les pages d'un même
			domaine fonctionnel. \\
			\hline
			\rowcolor{white}
			ORM & \textit{Object-Relational Mapping} : correspondance entre les tables
			d'une base de données et les objets du programme (ici SQLAlchemy). \\
			\hline
			\rowcolor{gray!8}
			Migration & Script versionné qui fait évoluer le schéma de la base de
			données sans perdre les données existantes. \\
			\hline
			\rowcolor{white}
			CRUD & \textit{Create, Read, Update, Delete} : les quatre opérations de
			base sur une donnée. \\
			\hline
			\rowcolor{gray!8}
			CSRF & \textit{Cross-Site Request Forgery} : attaque consistant à faire
			exécuter une action par un utilisateur authentifié à son insu. \\
			\hline
			\rowcolor{white}
			Git & Système de gestion de versions du code source. \\
			\hline
			\rowcolor{gray!8}
			Test unitaire & Programme qui vérifie automatiquement qu'une fonction
			produit bien le résultat attendu. \\
			\hline
			\rowcolor{white}
			WSGI & Interface standard entre un serveur web et une application Python. \\
			\hline
			\rowcolor{gray!8}
			ECC & École Centrale Casablanca. \\
			\hline
}

% ══════════════════════════════════════════════════════════════════════════════
\printstyledtoc
\pagenumbering{arabic}
\setcounter{page}{1}

% ══════════════════════════════════════════════════════════════════════════════
\mainsectionstar{Introduction}

Un investissement immobilier se juge sur des chiffres, mais ces chiffres ne
disent rien tant qu'on ignore ce que le client attend de son argent, et à quelle
échéance. C'est la leçon centrale des six semaines que j'ai passées chez
\entreprise, cabinet de conseil en investissement immobilier installé à
\villestage, du \datedebut{} au \datefin.

La mission qui nous a été confiée, à \binome{} et à moi-même, était en apparence
simple : doter le cabinet d'un outil informatique capable de calculer la
rentabilité d'un investissement immobilier et de présenter le résultat à un
client en rendez-vous. Le cabinet réalisait ces analyses au cas par cas, sur
tableur, sans support homogène ni document remis en fin d'entretien. Il fallait
industrialiser ce travail sans le dénaturer.

Nous avons donc commencé par ce que notre formation nous avait appris à faire :
écrire les formules, poser les conventions de calcul, construire un moteur
financier, l'entourer de tests. Cette partie du travail a fonctionné. Mais elle
reposait sur une hypothèse que nous n'avions jamais vérifiée --- qu'un
investissement immobilier se juge d'abord au rendement locatif et que tout bien
produit un loyer. À mi-parcours, un entretien de cadrage avec le tuteur a montré
que les deux dossiers dont le cabinet est le plus fier ne comportaient
\emph{aucun} loyer, et qu'aucun seuil de rentabilité ne servait à trancher : ce
qui décide, c'est l'horizon du client.

Ce rapport raconte ce déplacement. Il ne s'agit pas seulement d'un correctif
technique : c'est le passage d'un outil pensé depuis la formule à un outil pensé
depuis le geste métier, et c'est l'apprentissage que je retiens de ce stage.

La première partie présente le cabinet, son secteur et les enjeux auxquels il
fait face, puis le contexte et le contenu de la mission. La seconde analyse le
déroulement du stage --- la manière dont le travail a été organisé, les
difficultés rencontrées et la façon dont elles ont été traitées --- avant
d'en tirer un bilan des apports et des manques. La conclusion revient sur ce que
ce stage change à mon projet professionnel.

% ══════════════════════════════════════════════════════════════════════════════
\mainsection{L'organisme d'accueil et la mission}

\subsection{Le secteur : le conseil en investissement immobilier au Maroc}

Le marché immobilier marocain, et singulièrement celui du Grand Casablanca,
présente une caractéristique qui explique l'existence même de cabinets comme
\entreprise{} : il est à la fois très actif et faiblement outillé pour
l'investisseur particulier. Les biens s'échangent, les prix évoluent, des
quartiers entiers se transforment --- mais l'information sur ce que rapporte
réellement un bien reste dispersée, détenue par les agences, les promoteurs et
les intermédiaires, et rarement mise en forme pour permettre une comparaison.

\figurecredit{0.92}{casablanca_vue_aerienne.jpg}
{Vue aérienne de Casablanca. La densification du tissu urbain et l'ouverture de
	nouveaux quartiers sont les moteurs de la valorisation foncière sur laquelle
	repose une partie du conseil du cabinet.}
{Photographie Jimmy Baikovicius, Wikimedia Commons, CC BY-SA 2.0.}

Dans ce contexte, le conseil en investissement immobilier occupe une position
d'intermédiation particulière. Il ne s'agit ni de l'agence immobilière, qui vend
un bien, ni du conseiller financier réglementé, qui gère un patrimoine
mobilier. Le cabinet de conseil se situe entre les deux : il part du besoin d'un
client, cherche le bien qui y répond, chiffre ce que l'opération lui coûtera et
lui rapportera, et l'accompagne jusqu'à la remise des clés.

Trois traits du marché local structurent ce métier et se sont retrouvés,
directement, dans les choix de conception de l'outil :

\begin{itemize}
	\item \textbf{Le foncier pèse autant que le locatif.} Une part significative
	des opérations conseillées ne produit aucun loyer : achat d'un terrain
	conservé plusieurs années, construction en vue d'une revente. La valeur y
	vient de la plus-value, pas du rendement.
	\item \textbf{L'achat sur plan est courant.} La VEFA décale la mise en
	exploitation du bien de plusieurs mois à plusieurs années, alors que les
	mensualités d'emprunt, elles, courent dès la signature.
	\item \textbf{L'accès au crédit dépend du profil.} La nationalité et la
	situation professionnelle du client conditionnent l'accès au financement
	bancaire local, ce qui en fait des informations de premier entretien et non
	des données administratives secondaires.
\end{itemize}

\subsection{\entreprise{} : activité, offre et structure}

\entreprise{} est un cabinet de conseil en investissement immobilier basé à
\villestage. \acompleter{date de création, forme juridique (SARL, SA…),
	effectif, chiffre d'affaires ou volume de dossiers annuel --- à demander au
	tuteur} Le cabinet accompagne principalement des investisseurs particuliers,
résidents comme non-résidents, sur l'ensemble du cycle d'une acquisition.

L'offre du cabinet couvre quatre familles de biens --- terrains, villas,
appartements, immeubles --- sur toute la gamme de standing, de l'économique au
luxe, et sur des biens neufs, anciens ou vendus sur plan. Cette amplitude n'est
pas un argument commercial : c'est une contrainte de conception. Un outil qui ne
saurait traiter qu'un appartement loué laisserait de côté une partie importante
de l'activité réelle.

\begin{encadre}
	\textbf{\color{eccblue}Ce que le cabinet vend, en une phrase.} Non pas un bien,
	mais un accompagnement : la capacité à transformer un besoin exprimé en une
	acquisition chiffrée, puis à conduire le dossier jusqu'à la livraison.
\end{encadre}

Le processus de travail du cabinet, tel qu'il nous a été décrit lors de
l'entretien de cadrage, se déroule en six étapes stables :

\begin{enumerate}
	\item \textbf{Recueil du besoin} --- un premier entretien où sont notés le
	type de bien recherché, le standing, la superficie souhaitée, la
	distribution intérieure pour un appartement (chambres, salles de bains,
	salons, étage, orientation), les commodités attendues (transports, écoles,
	commerces), la zone géographique, le budget prévu, le mode de financement,
	l'objectif poursuivi et l'horizon d'investissement.
	\item \textbf{Recherche} --- identification des biens correspondant aux critères.
	\item \textbf{Présentation et visites} --- les biens retenus sont présentés,
	puis visités.
	\item \textbf{Compromis} --- signature du compromis ou du contrat de vente.
	\item \textbf{Acte notarié} --- finalisation de l'achat devant notaire.
	\item \textbf{Livraison} --- remise du bien au client.
\end{enumerate}

À chacune de ces étapes, une information systématique est recueillie sur le
client lui-même : nom, situation professionnelle, nationalité et budget
disponible. Ces quatre données sont demandées avant toute proposition.

\subsection{Position sur le marché et enjeux actuels}

\acompleter{position concurrentielle du cabinet : nombre de concurrents
	identifiés sur le Grand Casablanca, éléments de différenciation revendiqués
	--- à préciser avec le tuteur}

De mon point de vue d'observateur externe, trois enjeux structurent l'activité du
cabinet et donnent son sens à la mission qui nous a été confiée.

\textbf{Premier enjeu : la traçabilité de l'argumentaire.} Le conseil du cabinet
repose sur des calculs faits au cas par cas, souvent sur tableur. Ce mode de
travail fonctionne tant que le conseiller qui a construit le tableau est présent
pour l'expliquer. Il devient fragile dès qu'il faut retrouver, six mois plus
tard, pourquoi tel chiffre a été retenu, ou dès qu'un autre conseiller reprend le
dossier. Un outil partagé, dont les conventions de calcul sont écrites et
imprimées sous chaque tableau, répond directement à ce risque.

\textbf{Deuxième enjeu : la crédibilité en rendez-vous.} Le moment décisif du
conseil est l'entretien. Un client qui repart avec un document propre, chiffré et
daté n'est pas dans la même disposition qu'un client à qui l'on a montré un
tableur à l'écran. Le cabinet a exprimé ce besoin de restitution --- un rapport
remis en fin de rendez-vous --- comme une attente prioritaire.

\textbf{Troisième enjeu : la diversité des opérations.} Le cabinet traite des
dossiers dont la logique économique diffère radicalement : un appartement loué,
un terrain porté dix ans, un immeuble construit pour être revendu en deux ans.
Un outil qui n'en couvrirait qu'une famille serait un outil de démonstration,
pas un outil de travail.

\subsection{La mission : contexte, contenu et environnement de travail}

\subsubsection{Contexte et objectif}

La mission consistait à concevoir et développer \textbf{RentImmo}, une
application web interne au cabinet, permettant à un conseiller de saisir un
bien, de construire plusieurs hypothèses de financement, de les comparer et de
remettre au client une synthèse imprimable. L'objectif chiffré fixé au cadrage
était que l'ensemble de ce parcours tienne en moins de dix minutes, durée d'une
séquence de rendez-vous.

Le projet a été mené en binôme avec \binome, sur six semaines, avec un rythme de
livraison hebdomadaire : chaque dimanche, un rapport d'avancement était transmis,
et le code de la semaine était versionné dans un dépôt Git.

\subsubsection{Équipe, service et environnement de travail}

\acompleter{préciser : service ou entité de rattachement au sein du cabinet,
	nombre de personnes y travaillant, rythme de présence (présentiel, hybride,
	distanciel), fréquence des points avec le tuteur}

L'organisation du travail a reposé sur une répartition claire au sein du binôme
et sur un rythme de rendez-vous avec le tuteur. Le travail de développement
lui-même a été largement autonome : nous définissions les objectifs de la semaine
à partir d'une feuille de route établie en semaine~1, et le tuteur intervenait
comme référent métier --- c'est-à-dire sur les questions de fond (que calcule-t-on,
et pourquoi), non sur les choix techniques.

Cette autonomie a été à la fois le principal atout et le principal risque du
stage. Atout, parce qu'elle a permis d'avancer vite et de livrer l'intégralité du
périmètre fonctionnel. Risque, parce qu'un outil développé loin de son
utilisateur dérive vers ce que le développeur trouve intéressant plutôt que vers
ce dont le métier a besoin --- ce qui, comme la seconde partie de ce rapport le
détaille, s'est effectivement produit.

\subsubsection{Les tâches confiées}

\begin{table}[H]
	\centering
	\small
	\setstretch{1.05}
	\renewcommand{\arraystretch}{1.40}
	\begin{tabular}{|>{\raggedright\arraybackslash}p{4.3cm}|>{\raggedright\arraybackslash}p{10.6cm}|}
		\hline
		\rowcolor{eccblue!15}
		\textbf{\color{eccblue}Tâche} & \textbf{\color{eccblue}Objectif assigné} \\
		\hline
		\rowcolor{white}
		Cadrage fonctionnel & Écrire le cahier des charges : périmètre, formules,
		conventions de calcul, modèle de données. \\
		\hline
		\rowcolor{gray!8}
		Moteur financier & Implémenter en Python les calculs d'acquisition, de
		crédit, de rendement, de VAN et de TRI, et les couvrir de tests. \\
		\hline
		\rowcolor{white}
		Application web & Développer les écrans de saisie et de restitution, avec
		authentification et cloisonnement des dossiers par conseiller. \\
		\hline
		\rowcolor{gray!8}
		Restitution client & Produire un rapport PDF et un classeur Excel aux
		couleurs du cabinet. \\
		\hline
		\rowcolor{white}
		Documentation & Rédiger un guide d'utilisation, un script de démonstration
		et les rapports hebdomadaires. \\
		\hline
		\rowcolor{gray!8}
		Mise en ligne & Rendre l'application accessible en ligne et documenter la
		procédure de déploiement. \\
		\hline
	\end{tabular}
\end{table}

Deux tâches ne figuraient pas dans la commande initiale et se sont imposées en
cours de stage : la \textbf{conduite d'un entretien de cadrage métier} avec le
tuteur, sous forme d'un questionnaire écrit de huit questions, et la
\textbf{refonte de l'application} qui a suivi les réponses reçues --- refonte
dont la dernière étape, l'automatisation de la construction et de la
confrontation des montages de financement, a été menée en toute fin de stage.
Ces deux tâches constituent, avec le recul, le cœur de ce que le stage m'a
appris.

% ══════════════════════════════════════════════════════════════════════════════
\mainsection{Analyse du déroulement et des apports du stage}

\subsection{L'organisation du travail : une règle et un rythme}

Une règle a été posée dès la première semaine et a structuré tout le
projet : \textbf{aucune ligne de logique financière ne serait écrite avant que
	les conventions de calcul ne soient rédigées}. Elle peut sembler formaliste ;
elle s'est révélée décisive. La plupart des arbitrages délicats du projet
n'étaient pas des problèmes de code mais des questions de définition : que
compte-t-on exactement dans les charges d'un rendement « net » ? Les travaux
font-ils partie du coût d'acquisition ou d'une dépense d'exploitation ?
L'assurance emprunteur se calcule-t-elle sur le capital initial ou sur le capital
restant dû ? Un champ laissé vide vaut-il zéro ou « hériter de la valeur par
défaut » ?

Chacune de ces questions admet plusieurs réponses défendables. Aucune n'admet
deux réponses simultanées dans un même programme. Les avoir tranchées et écrites
avant de coder a évité les incohérences ; les avoir affichées dans l'interface et
imprimées dans les exports a donné au conseiller de quoi défendre chaque chiffre
devant son client.

\begin{table}[H]
	\centering
	\small
	\setstretch{1.05}
	\renewcommand{\arraystretch}{1.40}
	\begin{tabular}{|>{\centering\arraybackslash}p{1.2cm}|>{\raggedright\arraybackslash}p{4.6cm}|>{\raggedright\arraybackslash}p{8.6cm}|}
		\hline
		\rowcolor{eccblue!15}
		\textbf{\color{eccblue}Sem.} & \textbf{\color{eccblue}Objet} &
		\textbf{\color{eccblue}Livré} \\
		\hline
		\rowcolor{white}
		1 & Cadrage & Cahier des charges, feuille de route, conventions de calcul,
		modèle de données. \\
		\hline
		\rowcolor{gray!8}
		2 & Moteur et socle & Moteur financier en Python pur, modèles de données,
		authentification, parcours client $\rightarrow$ projet $\rightarrow$ scénario. \\
		\hline
		\rowcolor{white}
		3 & Restitution & Graphiques, comparaison de scénarios côte à côte,
		aperçu recalculé à chaque saisie. \\
		\hline
		\rowcolor{gray!8}
		4 & Documents client & Travaux détaillés par postes, export PDF, export
		Excel, duplication de scénario. \\
		\hline
		\rowcolor{white}
		5 & Consolidation & Cas réels vérifiés à la main, cas limites, jeu de
		démonstration, guide d'utilisation. \\
		\hline
		\rowcolor{gray!8}
		6 & Cadrage métier et refonte & Entretien avec le tuteur, refonte du
		modèle et de l'interface, mise en ligne, remise. \\
		\hline
	\end{tabular}
	\caption{Déroulé effectif du stage, semaine par semaine.}
\end{table}

Le rythme hebdomadaire --- un rapport chaque dimanche, un jeu de tests au vert à
chaque fin de semaine --- a joué un rôle qui dépasse le suivi administratif. Il
a imposé de terminer ce qui était commencé avant d'ouvrir un nouveau chantier.
À la fin de la semaine~5, le développement prévu était achevé sans dette
technique, ce qui a rendu possible la refonte de la semaine~6. Sans cette
avance, l'entretien de cadrage n'aurait débouché sur rien d'autre qu'une liste
de regrets.

\subsection{Le tournant du stage : confronter l'outil au métier réel}

\subsubsection{Une hypothèse jamais vérifiée}

À la fin de la semaine~5, l'application faisait exactement ce que le cahier des
charges annonçait. Elle calculait un coût d'acquisition, une mensualité, trois
rendements, un cash-flow, une VAN et un TRI ; elle affichait le TRI et la VAN en
tête de la page de résultats, parce que ce sont les indicateurs les plus
sophistiqués ; elle supposait qu'un bien produit un loyer, parce que c'est ce
qu'un bien fait dans les manuels.

Nous avons alors adressé au tuteur un questionnaire de huit questions sur sa
méthode réelle : quels indicateurs regarde-t-il en premier ? Existe-t-il des
seuils de rentabilité ? Quelles informations recueille-t-il au premier
entretien ? Peut-il décrire deux dossiers récents ?

Les réponses ont invalidé trois hypothèses de conception sur quatre.

\begin{citationbox}
	« Les deux critères les plus récurrents sont le \textbf{rendement net} et le
	\textbf{cash-flow} que le bien peut générer par la suite. Viennent ensuite les
	autres éléments cités. »

	\vspace{0.4em}
	« Un investissement est considéré comme bon lorsqu'il permet au client de
	générer de la valeur. Tout dépend donc de son \textbf{horizon
		d'investissement} et du \textbf{facteur temps}. »

	\vspace{0.4em}
	{\footnotesize\color{smoke} Extraits de l'entretien de cadrage métier,
		août 2026.}
\end{citationbox}

Trois enseignements en découlent. D'abord, le TRI et la VAN --- les indicateurs
sur lesquels nous avions le plus travaillé --- sont des arguments, pas des
critères de décision : le conseiller regarde d'abord le rendement net et le
cash-flow. Ensuite, \textbf{il n'existe aucun seuil de rentabilité} : « un
investissement peut être rentable pour une personne et ne pas l'être pour une
autre », ce qui interdit tout verdict automatique du type « bonne affaire / à
éviter ». Enfin, et c'est le point le plus lourd de conséquences, les deux
dossiers cités en exemple par le tuteur ne comportaient \emph{aucun loyer}.

\subsubsection{Les deux dossiers de référence}

\begin{table}[H]
	\centering
	\small
	\setstretch{1.05}
	\renewcommand{\arraystretch}{1.40}
	\begin{tabular}{|>{\raggedright\arraybackslash}p{2.6cm}|>{\raggedright\arraybackslash}p{5.9cm}|>{\raggedright\arraybackslash}p{5.9cm}|}
		\hline
		\rowcolor{eccblue!15}
		& \textbf{\color{eccblue}Client 1 --- portage foncier} &
		\textbf{\color{eccblue}Client 2 --- construction-revente} \\
		\hline
		\rowcolor{white}
		Budget & environ 1 000 000 MAD & environ 1 000 000 MAD \\
		\hline
		\rowcolor{gray!8}
		Bien & Terrain de 3 ha en zone ciblée & Terrain viabilisé de 500 m²
		à 1 100 000 MAD \\
		\hline
		\rowcolor{white}
		Opération & Conservé, sans mise en valeur & Immeuble R+4, 20 appartements,
		revendus \\
		\hline
		\rowcolor{gray!8}
		Horizon & 10 ans & 2 ans \\
		\hline
		\rowcolor{white}
		Résultat & Un projet d'État s'implante à 1 km quatre ans plus tard ;
		morcellement et revente de 2 ha, 1 ha conservé en réserve $\rightarrow$
		\textbf{environ 15 000 000 MAD} de bénéfice &
		\textbf{5 000 000 MAD} de plus-value nette \\
		\hline
	\end{tabular}
	\caption{Les deux dossiers de référence décrits par le tuteur.}
\end{table}

Ces deux dossiers disent, à eux seuls, ce que l'outil devait devenir. À la sortie
du client~2, celui-ci semble avoir réalisé la meilleure opération : cinq millions
en deux ans. Sur l'horizon du client~1, le classement s'inverse. Le même capital,
le même cabinet, deux résultats qui ne se comparent qu'à condition de fixer
d'abord la durée. C'est exactement ce que le tuteur appelle le « facteur temps ».

\figurecredit{0.86}{casablanca_marina_chantier.jpg}
{Un chantier en cours à Casablanca. Pour un bien acheté sur plan, les mensualités
	d'emprunt courent pendant que le bien ne produit encore rien : c'est ce
	décalage que le moteur de calcul a dû apprendre à représenter.}
{Photographie Farid Mernissi, Wikimedia Commons, CC BY-SA 4.0.}

\subsubsection{Ce qui a été refait}

La semaine~6 a été consacrée à la refonte. Elle a porté sur trois niveaux.

\textbf{Le modèle de données.} La fiche client a reçu les trois informations que
le cabinet recueille systématiquement et que nous ignorions : situation
professionnelle, nationalité, budget disponible. Un \textbf{brief de recherche} a
été ajouté, reprenant tous les critères du premier entretien, ainsi que
l'objectif du client et son horizon. Le dossier a reçu un type d'opération
(locatif ou terrain / revente), une étape d'avancement suivant les six phases du
métier, et un délai de livraison en mois.

\textbf{Le moteur de calcul.} Une opération sans loyer est désormais traitée
nativement : toute la valeur y vient de la plus-value. Le délai de livraison
diffère la mise en exploitation --- tant que le bien n'est pas livré il ne
produit ni loyer ni charges, alors que les annuités d'emprunt courent déjà --- et
une livraison en cours d'année est proratisée. Le cash-flow mis en avant est
celui d'une année pleinement exploitée, et non celui de l'année~1 qui ne
représente rien pour un bien encore en chantier.

\textbf{La restitution.} L'ordre des indicateurs a été inversé sur tous les
supports : rendement net et cash-flow en tête, valeur créée ensuite, TRI et VAN
en second rang. Tout verdict automatique a été supprimé ; à la place, l'objectif
et l'horizon du client sont rappelés à côté des chiffres. L'étude automatique
ajoutée en fin de stage (\S\,\ref{sec:saisie}) ne revient pas sur ce principe : elle
\emph{propose} un montage de financement pour un bien donné, jamais un jugement
sur l'investissement lui-même, et toujours relativement à l'objectif que le
client a déclaré. Les écrans s'adaptent
au type d'opération : pas de rendement locatif affiché pour un terrain, la
plus-value prend sa place.

\begin{encadrecle}
	\textbf{Validation.} Les deux dossiers du tuteur sont devenus des
	\textbf{tests automatisés}, vérifiables à la main. Client~1 : 1 000 000 + 7 \%
	de frais = 1 070 000 investis, revente à 16 070 000, soit
	\textbf{15 000 000 de valeur créée} et un TRI d'environ 31 \% par an.
	Client~2 : 1 100 000 + 7 \% + 8 000 000 de construction = 9 177 000, revente à
	14 177 000, soit \textbf{5 000 000 de valeur créée}. Un troisième test vérifie
	que \textbf{le classement des deux dossiers s'inverse avec l'horizon} : la
	conclusion métier du tuteur est ainsi inscrite dans le code, et toute
	régression future la ferait échouer.
\end{encadrecle}

\subsection{Difficultés rencontrées et manière dont elles ont été traitées}

\subsubsection{Un algorithme qui refuse de converger}

Le taux de rentabilité interne n'a pas de formule fermée : il se cherche par
approximations successives. La méthode de Newton-Raphson, rapide, peut diverger
ou sortir du domaine admissible lorsque les flux changent plusieurs fois de
signe --- exactement le cas d'un investissement immobilier à crédit, où le
cash-flow est négatif pendant des années avant l'encaissement de la revente.

La solution retenue combine deux méthodes : Newton-Raphson d'abord, et, en cas
d'échec, un repli sur une \textbf{bissection bornée} à un intervalle réaliste. Et
lorsque aucune solution n'existe, l'outil affiche explicitement l'absence de
valeur \textbf{plutôt qu'un chiffre faux}. Ce dernier point est un choix de
conception, non une limite technique : un chiffre trompeur affiché en rendez-vous
est plus dommageable qu'une case vide.

\subsubsection{Un test qui échoue parce que le vérificateur se trompe}

En semaine~5, un test construit sur un cas réel a échoué. Le réflexe a été de
chercher l'erreur dans le moteur. Elle n'y était pas : la valeur de référence,
que j'avais calculée à la main, était fausse de quelques dirhams sur une
multiplication. Le programme, lui, avait raison.

L'anecdote est mineure ; l'enseignement ne l'est pas. Le principe posé au
cadrage --- \textbf{chaque chiffre vérifié deux fois, par le code et par un
	calcul indépendant} --- avait été conçu pour détecter les erreurs du programme.
Il a détecté une erreur de l'humain. C'est le même dispositif qui protège dans
les deux sens, et j'ai cessé, ce jour-là, de considérer les tests comme une
formalité de fin de développement.

\subsubsection{Un bug invisible en développement, fatal en production}

Le bug le plus instructif du stage n'a jamais provoqué d'erreur sur nos postes.
La configuration de l'application lisait ses paramètres --- dont la clé secrète
et l'adresse de la base de données --- depuis un fichier d'environnement, mais
ce fichier était chargé \emph{après} que la configuration ait été évaluée. En
développement, l'outil en ligne de commande de Flask charge ce fichier de son
côté : le défaut restait masqué. Sous le serveur de production, ce filet n'existe
pas : l'application serait partie avec ses valeurs par défaut, c'est-à-dire une
clé secrète publique et une base de données vide.

Le correctif tient en une ligne déplacée. Le trouver a demandé de comprendre
\textbf{dans quel ordre les choses se produisent au démarrage}, et de vérifier
l'application dans les deux modes d'exécution au lieu d'un seul. C'est le type de
défaut qu'aucun test fonctionnel n'attrape, parce que l'environnement de test
ressemble à l'environnement de développement.

\subsubsection{Trop de champs à remplir}
\label{sec:saisie}

Une difficulté de conception, celle-là. À force d'ajouter les critères
recueillis par le cabinet, les formulaires étaient devenus trop longs pour être
remplis en rendez-vous --- ce qui contredisait l'exigence des dix minutes fixée
au cadrage.

La première réponse a consisté à distinguer trois catégories de champs plutôt
qu'à en supprimer arbitrairement : ce qui est \textbf{obligatoire} parce que sans
lui le calcul n'a pas de sens ou parce que le cabinet le recueille
systématiquement ; ce qui est \textbf{suggéré} par des listes de valeurs
fréquentes que le conseiller peut ignorer ; et ce qui est \textbf{replié}
derrière un bloc dépliable, visible seulement quand on en a besoin. Ce travail a
fait l'objet d'une note écrite justifiant chaque arbitrage, précisément parce
qu'un choix d'interface qu'on ne sait pas expliquer est un choix qu'on n'aurait
pas dû faire.

Cette réponse traitait le symptôme. La cause était ailleurs, et je ne l'ai
comprise qu'en regardant ce que le conseiller faisait réellement de l'écran de
scénario : il y saisissait onze valeurs --- apport, taux, assurance, durée,
horizon, revalorisations, frais de revente, taux d'actualisation --- pour
découvrir que le montage ne convenait pas, puis recommençait. \textbf{L'outil
lui demandait de deviner la réponse avant de la calculer.}

D'où la dernière fonctionnalité développée, et celle dont je suis le plus
satisfait : l'\textbf{étude automatique}. Le conseiller décrit le bien et le
client ; l'application construit elle-même la famille des montages plausibles
--- paiement comptant, crédit à apport minimal, renforcé ou maximal, sur quinze,
vingt et vingt-cinq ans --- les calcule tous, les classe, et \textbf{propose
celui qui répond le mieux à l'objectif déclaré du client}, en français courant :
« avec un apport de tant, le bien paie ses charges et laisse tant par mois ; au
bout de vingt ans, l'opération aura rapporté tant ».

La difficulté n'était pas de programmer ce classement, mais de ne pas trahir au
passage ce que le cabinet nous avait dit --- \emph{« un investissement peut être
rentable pour une personne et ne pas l'être pour une autre »}. Proposer un
montage sans réintroduire un seuil de rentabilité universel demandait quatre
garde-fous, tous visibles à l'écran :

\begin{itemize}
	\item le classement est \textbf{relatif} : chaque critère est ramené entre le
	meilleur et le moins bon des montages \emph{de cette étude}, pour un même bien
	et un même horizon. Un score ne vaut donc qu'à l'intérieur d'une page, et deux
	études ne se comparent pas ;
	\item la \textbf{pondération est affichée} --- le cash-flow pèse 40~\% pour un
	objectif de revenu locatif, la valeur créée 40~\% pour un objectif de
	plus-value : le conseiller voit sur quoi l'outil a tranché, et peut le
	contester ;
	\item l'écran dit \textbf{ce que les autres montages font mieux}. C'est
	typiquement là qu'apparaît l'effet de levier : l'achat comptant l'emporte
	souvent sur le cash-flow tout en rendant moins par dirham engagé. Un
	argumentaire qui ne nommerait pas ce que l'on perd serait de la vente, pas du
	conseil ;
	\item chaque montage affiche sa \textbf{composition complète}, paramètre par
	paramètre, avec l'origine de chaque valeur : dossier, brief du client, zone de
	marché, calcul, ou hypothèse par défaut. Un chiffre qu'on ne peut pas
	justifier devant un client n'a pas sa place dans une proposition.
\end{itemize}

Un test résume l'intention mieux qu'un paragraphe : \emph{à bien identique et
budget identique}, l'outil propose le crédit sur vingt-cinq ans à un client qui
cherche du revenu, et celui sur quinze ans à un client qui vise la plus-value.
La position du cabinet est ainsi devenue une propriété vérifiée du programme.

\subsubsection{Deux historiques de code sans ancêtre commun}

Au moment de publier le projet, le dépôt distant et le dépôt local se sont
révélés avoir été créés indépendamment : deux historiques Git sans aucun commit
partagé, que l'outil refusait donc de fusionner. La tentation était de tout
écraser. J'ai préféré vérifier d'abord, fichier par fichier, qu'aucun travail de
mon binôme ne serait perdu, puis fusionner explicitement les deux histoires en
donnant la priorité à la version locale --- vérifiée --- des fichiers en conflit.

Une seconde opération a suivi : une image de référence avait été versionnée par
inadvertance. La retirer du suivi ne suffit pas, car elle reste dans l'historique
d'un dépôt public. Il a fallu réécrire cet historique et republier. J'ai retenu
de cet épisode que \textbf{ce qui est publié est difficile à dépublier}, et que
la prudence se joue avant l'envoi, pas après.

\subsubsection{Des graphiques qui ne trompent pas}

Dernière difficulté, moins technique et plus inattendue : la rigueur graphique.
Un outil destiné à convaincre en rendez-vous ne peut pas se permettre des
graphiques ambigus. Des règles strictes ont été appliquées --- un seul axe
vertical par graphique, couleurs attribuées dans un ordre fixe et contrôlées
pour les principales formes de daltonisme, légende systématique dès deux séries,
tableau chiffré en regard de chaque figure --- quitte à écarter des
représentations séduisantes mais trompeuses, comme les doubles axes ou les
jauges. Aucune information ne repose sur la couleur seule.

\subsection{Le résultat livré}

L'application livrée couvre l'intégralité du périmètre du cahier des charges,
révisé après l'entretien de cadrage, et s'appuie sur une suite de
\textbf{96 tests automatisés} au vert à la remise. Le parcours nominal tient en
trois écrans --- fiche client et brief, dossier du bien, étude automatique ---
et la saisie d'un montage à la main, qui était le chemin obligé, n'est plus qu'un
recours pour un financement déjà négocié.

\figurecredit{0.96}{app_resultats_haut.png}
{Page de résultats d'un scénario, après refonte : le rendement net et le
	cash-flow occupent la première ligne, la valeur créée, le TRI et la VAN la
	seconde. Cet écran n'affiche aucun verdict ; l'objectif et l'horizon du client
	sont rappelés en tête.}
{Capture d'écran de l'application développée pendant le stage.}

\vspace{-0.4em}
\begin{table}[H]
	\centering
	\small
	\setstretch{1.05}
	\renewcommand{\arraystretch}{1.40}
	\begin{tabular}{|>{\raggedright\arraybackslash}p{4.5cm}|>{\raggedright\arraybackslash}p{10.4cm}|}
		\hline
		\rowcolor{eccblue!15}
		\textbf{\color{eccblue}Livrable} & \textbf{\color{eccblue}État à la remise} \\
		\hline
		\rowcolor{white}
		Application web & Parcours complet, de la fiche client au document remis,
		déployable en ligne. \\
		\hline
		\rowcolor{gray!8}
		Étude automatique & Construction, confrontation et proposition motivée des
		montages de financement, en langage clair, sans saisie supplémentaire. \\
		\hline
		\rowcolor{white}
		Moteur financier & Module Python indépendant, testé contre des valeurs de
		référence externes et contre les deux dossiers du cabinet. \\
		\hline
		\rowcolor{gray!8}
		Cahier des charges & Rédigé en semaine~1, complété d'une note d'écart
		après l'entretien de cadrage. \\
		\hline
		\rowcolor{white}
		Documentation & Guide d'utilisation, script de démonstration minuté,
		procédure de déploiement, note de choix d'interface. \\
		\hline
		\rowcolor{gray!8}
		Rapports & Cinq rapports hebdomadaires et un rapport final de projet. \\
		\hline
	\end{tabular}
	\caption{Livrables remis au cabinet le \datefin.}
\end{table}

\subsection{Analyse des apports du stage}

\subsubsection{Points forts}

\textbf{Un projet complet, de bout en bout.} J'ai vu un même travail passer de la
question métier au document remis à un client. C'est rare dans un cadre
scolaire, où l'on hérite le plus souvent d'un énoncé déjà formulé. Ici, une part
substantielle de la difficulté résidait dans la formulation elle-même.

\textbf{Un utilisateur réel, avec une opinion.} L'entretien de cadrage a produit
plus de valeur en huit questions que trois semaines de développement autonome.
J'ai appris que l'accès à l'utilisateur n'est pas une étape du projet mais sa
condition, et qu'il faut aller la chercher tôt plutôt que la subir tard.

\textbf{La contrainte de l'explicabilité.} Chaque chiffre devait pouvoir être
justifié en une phrase devant un client. Cette contrainte a écarté d'elle-même
les solutions élégantes mais opaques et m'a habitué à un critère que je n'avais
pas : \emph{si je ne sais pas l'expliquer, je ne dois pas le livrer}.

\subsubsection{Points faibles}

\textbf{L'entretien de cadrage est arrivé trop tard.} Il aurait dû être conduit
en semaine~1, pas en semaine~6. Cinq semaines de travail ont reposé sur des
hypothèses non vérifiées ; la refonte a été possible parce que l'architecture
s'y prêtait et parce que nous avions de l'avance, mais c'est une chance, pas une
méthode. Rétrospectivement, la faute est claire : nous avons pris le silence de
départ pour un accord.

\textbf{Aucun test en situation réelle.} L'application n'a pas été utilisée
devant un vrai client au cours du stage. Le jeu de démonstration et le script de
rendez-vous en tiennent lieu, mais ils restent des simulations construites par
ceux-là mêmes qui ont écrit l'outil.

\textbf{Des zones métier laissées en suspens.} L'échéancier des paiements en
VEFA, le calcul de la capacité d'emprunt et le traitement fiscal de la plus-value
n'ont pas été traités. Ils sont documentés comme tels, ce qui est honnête, mais
ce sont des manques.

\subsubsection{Compétences acquises}

\begin{table}[H]
	\centering
	\small
	\setstretch{1.05}
	\renewcommand{\arraystretch}{1.40}
	\begin{tabular}{|>{\raggedright\arraybackslash}p{4.2cm}|>{\raggedright\arraybackslash}p{10.7cm}|}
		\hline
		\rowcolor{eccblue!15}
		\textbf{\color{eccblue}Domaine} & \textbf{\color{eccblue}Ce que j'ai appris à faire} \\
		\hline
		\rowcolor{white}
		Techniques & Structurer une application web (Flask, base de données,
		migrations), isoler une logique métier pure et la tester, générer des
		documents PDF et Excel, déployer une application en ligne. \\
		\hline
		\rowcolor{gray!8}
		Financières & Manipuler les indicateurs d'un investissement --- rendement,
		cash-flow, VAN, TRI, amortissement d'un emprunt --- et surtout comprendre
		ce que chacun ne dit pas. \\
		\hline
		\rowcolor{white}
		Méthodologiques & Écrire les conventions avant le code, vérifier chaque
		chiffre par deux voies indépendantes, versionner et documenter au rythme
		hebdomadaire. \\
		\hline
		\rowcolor{gray!8}
		Relationnelles & Formuler des questions métier exploitables, entendre une
		réponse qui invalide son propre travail, et transformer un retour critique
		en plan d'action plutôt qu'en justification. \\
		\hline
	\end{tabular}
\end{table}

\subsubsection{Compétences qui m'ont manqué}

Trois manques se sont fait sentir. D'abord, une \textbf{culture du marché
	immobilier local} : je découvrais en même temps le métier et son vocabulaire, ce
qui m'a fait perdre du temps sur des notions qu'un entretien préalable aurait
levées en une heure. Ensuite, des \textbf{méthodes de recueil du besoin} : je ne
disposais d'aucune technique structurée d'entretien utilisateur, et le
questionnaire que nous avons envoyé au tuteur a été improvisé --- il aurait été
plus efficace, et surtout plus précoce, si j'avais su comment on conduit
ce type d'échange. Enfin, des \textbf{réflexes d'exploitation} : la différence
entre « ça marche sur ma machine » et « ça marche en production » ne m'était pas
familière, et c'est un bug de configuration qui me l'a apprise.

\subsubsection{Lien avec la formation à l'ECC}

Plusieurs acquis de la première année ont été directement mobilisés.
La \textbf{programmation} et l'\textbf{algorithmique} ont servi au-delà de la
syntaxe : la recherche du TRI est une résolution numérique d'équation, et le
choix Newton-Raphson puis bissection est un arbitrage classique entre vitesse de
convergence et robustesse --- exactement le type de raisonnement vu en cours de
méthodes numériques. Les \textbf{mathématiques financières} ont fourni
l'actualisation et le calcul d'annuités constantes. La \textbf{communication
	écrite} a été sollicitée en permanence : cahier des charges, rapports
hebdomadaires, guide d'utilisation.

\acompleter{citer ici les intitulés exacts des cours de 1A concernés
	--- l'exercice attend des références précises à la formation}

\subsection{Enseignements sur le fonctionnement de l'entreprise}

Trois observations sur l'organisation du cabinet me paraissent mériter d'être
consignées.

\textbf{La décision appartient au client, et cela structure tout l'outillage.}
Ma première réaction, à la lecture de la réponse « il n'existe pas de seuil »,
a été de la trouver imprécise. Elle est exactement l'inverse : le cabinet a
choisi de ne pas se substituer au client, et cette position se traduit
directement dans le produit --- pas de seuil de rentabilité, mais des chiffres
présentés à côté de l'objectif et de l'horizon que le client a lui-même
exprimés. Une règle de conduite professionnelle est ici devenue une contrainte
technique : c'est elle qui a dicté, jusque dans le détail, la manière dont
l'étude automatique classe les montages --- relativement, pour un client donné,
pondération affichée, et en disant ce que les montages écartés faisaient mieux.

\textbf{Le savoir-faire est dans les têtes, pas dans les procédures.} Le
processus en six étapes n'était écrit nulle part avant que nous ne le
transcrivions : il existait comme une pratique partagée. Un cabinet de cette
taille fonctionne très bien ainsi tant que ses effectifs sont stables, et la
formalisation --- que notre outil impose de fait --- est autant un risque
(rigidifier ce qui était souple) qu'un progrès.

\textbf{La communication passe par l'écrit quand elle doit engager.} Les
échanges informels ont porté l'essentiel du quotidien, mais les décisions qui
comptaient --- les conventions de calcul, les réponses au questionnaire, les
écarts constatés --- ont toutes été fixées par écrit. C'est ce qui a permis, en
semaine~6, de mesurer précisément l'écart entre ce que nous avions construit et
ce que le métier attendait : les deux étaient documentés.

% ══════════════════════════════════════════════════════════════════════════════
\mainsectionstar{Conclusion}

Six semaines après, l'outil livré fait autre chose que ce que nous avions prévu
d'écrire. Il ne décrète pas qu'un investissement est bon, et il ne suppose plus
qu'un bien rapporte un loyer. Il présente à un conseiller ce que son client a dit
vouloir, sur la durée que son client a fixée, et lui donne de quoi le justifier
chiffre par chiffre. Quand il finit par proposer quelque chose --- le montage de
financement le mieux placé parmi ceux qu'il a lui-même construits --- il affiche
en même temps la règle qui l'a fait choisir et ce que les autres options
faisaient mieux, de sorte que la proposition reste discutable. Ce déplacement --- d'un outil qui décide
vers un outil qui documente une décision --- est le résultat le plus solide du
stage, et il ne vient pas d'une compétence technique : il vient d'avoir posé des
questions à quelqu'un qui exerce le métier.

Sur le plan technique, ce stage m'a confirmé un goût. J'ai aimé construire un
système dont chaque partie est vérifiable, et j'ai découvert que la partie la
plus intéressante n'était pas le code mais l'endroit où le code rencontre une
règle de métier. Un moteur financier qui donne le bon chiffre est une affaire de
quelques centaines de lignes ; savoir \emph{quel} chiffre doit être donné est
une affaire d'entretien, de convention écrite et de restitution.

C'est ce qui oriente la suite. Pour mon prochain stage, je chercherai un
environnement où cette rencontre est au centre : ingénierie logicielle appliquée
à un domaine métier exigeant --- finance, énergie, industrie --- plutôt que
développement pour lui-même. Je souhaite aussi être confronté à une équipe plus
étoffée : ce stage s'est déroulé en binôme, avec une autonomie large, et je n'ai
donc pas éprouvé ce que signifie tenir sa place dans une organisation de
plusieurs dizaines de personnes, avec ses processus de revue et ses arbitrages
collectifs.

Enfin, deux lacunes identifiées appellent un plan concret. La \textbf{conduite
	d'entretien utilisateur} : je compte m'y former par la lecture méthodique et
en imposant, dès le premier jour de ma prochaine mission, un entretien de
cadrage écrit avec l'utilisateur final --- l'erreur commise ici est trop coûteuse
pour être répétée. La \textbf{mise en exploitation} : je veux comprendre ce qui
se passe entre un programme qui fonctionne et un service qui tourne, en
pratiquant sur mes propres projets le déploiement, la configuration et la
supervision, plutôt qu'en les découvrant à la veille d'une remise.

\begin{encadrecle}
	\textbf{En une phrase.} J'ai commencé ce stage en croyant qu'il s'agissait de
	calculer juste ; je le termine en sachant qu'il s'agissait d'abord de calculer
	\emph{ce qu'il fallait}, et que cette question-là ne se résout pas devant un
	écran.
\end{encadrecle}

% ══════════════════════════════════════════════════════════════════════════════
\mainsectionstar{Sources}

Conformément aux consignes de rédaction, les éléments repris de sources externes
sont signalés ci-dessous. Le contenu métier de ce rapport provient d'une source
primaire unique : l'entretien de cadrage conduit avec le tuteur.

\begin{itemize}
	\item \textbf{Entretien de cadrage métier}, \entreprise, août 2026 ---
	questionnaire écrit de huit questions et réponses du tuteur. Les extraits
	cités sont reproduits littéralement ; l'intégralité est consignée dans la
	note d'écart remise au cabinet.
	\item \textbf{Documents produits pendant le stage} : cahier des charges
	(semaine~1), rapports hebdomadaires n\textsuperscript{os}~1 à 5, rapport final
	de projet, guide d'utilisation, note de choix d'interface, note d'écart et
	d'adaptation. Ces documents sont les nôtres.
	\item \textbf{Documentations techniques officielles} consultées pour la mise
	en œuvre : Flask, SQLAlchemy, WTForms, WeasyPrint, openpyxl, Chart.js,
	pytest, Git.
	\item \textbf{Photographies} : Wikimedia Commons, sous licences Creative
	Commons ; auteur et licence indiqués sous chaque figure.
	\item \textbf{Captures d'écran} : productions du stage, réalisées sur
	l'application développée avec son jeu de données de démonstration.
\end{itemize}

% ══════════════════════════════════════════════════════════════════════════════
\newpage
\appendix
\mainsectionstar{Annexes}
\renewcommand{\thesubsection}{\Alph{subsection}}
\setcounter{subsection}{0}

\subsection{Le parcours de l'application, écran par écran}

\figurecredit{0.94}{app_vitrine.png}
{Page d'accueil publique : le cabinet, sa méthode en six étapes et ses questions
	fréquentes. C'est la seule page visible sans compte.}{}

\figurecredit{0.94}{app_tableau_de_bord.png}
{Tableau de bord du conseiller : l'avancement des dossiers suit les six étapes du
	processus décrit par le cabinet.}{}

\figurecredit{0.94}{app_brief.png}
{Brief de recherche : les critères recueillis au premier entretien. Ce formulaire
	n'existait pas avant l'entretien de cadrage.}{}

\newpage
\figurecredit{0.94}{app_projet_terrain.png}
{Dossier d'une opération sans loyer --- un terrain porté sur dix ans. L'écran
	s'adapte : aucun rendement locatif n'est demandé ni affiché.}{}

\figurecredit{0.94}{app_comparaison.png}
{Comparaison de trois montages de financement pour un même bien : l'écran
	d'arbitrage du rendez-vous.}{}

\newpage
\figurecredit{0.94}{app_etude_carte.png}
{Étude automatique : les sept montages possibles ont été construits, calculés et
	classés pour l'objectif déclaré du client ; le mieux placé est proposé en une
	phrase, avec ses quatre chiffres, ce qui plaide pour lui, ce que les autres
	font mieux, et ce qu'il faut dire au client avant qu'il s'engage.}{}

\subsection{Le processus métier du cabinet}

\begin{center}
	\begin{tikzpicture}[node distance=0.42cm, font=\small]
		\tikzset{etape/.style={rectangle, rounded corners=3pt, draw=eccblue!60,
				fill=eccblue!8, text width=2.35cm, align=center, minimum height=1.5cm,
				inner sep=4pt}}
		\node[etape] (a) {\textbf{1}\\Recueil du besoin};
		\node[etape, right=of a] (b) {\textbf{2}\\Recherche};
		\node[etape, right=of b] (c) {\textbf{3}\\Présentation et visites};
		\node[etape, right=of c] (d) {\textbf{4}\\Compromis};
		\node[etape, below=0.9cm of b] (e) {\textbf{5}\\Acte notarié};
		\node[etape, right=of e] (f) {\textbf{6}\\Livraison};
		\draw[-{Stealth[length=5pt]}, eccblue!70, thick] (a) -- (b);
		\draw[-{Stealth[length=5pt]}, eccblue!70, thick] (b) -- (c);
		\draw[-{Stealth[length=5pt]}, eccblue!70, thick] (c) -- (d);
		\draw[-{Stealth[length=5pt]}, eccblue!70, thick] (d.south) |- ([yshift=0.45cm]e.north) -- (e.north);
		\draw[-{Stealth[length=5pt]}, eccblue!70, thick] (e) -- (f);
		\node[below=0.35cm of e, text width=12cm, align=center,
		font=\footnotesize\itshape, color=darkgray]
		at ([xshift=1.6cm,yshift=-0.9cm]e.south)
		{Les six étapes suivies par le cabinet, transcrites depuis l'entretien de
			cadrage. Elles sont devenues les six statuts d'un dossier dans
			l'application.};
	\end{tikzpicture}
\end{center}

\subsection{Ce qui a changé après l'entretien de cadrage}

\begin{table}[H]
	\centering
	\small
	\setstretch{1.05}
	\renewcommand{\arraystretch}{1.40}
	\begin{tabular}{|>{\centering\arraybackslash}p{0.7cm}|>{\raggedright\arraybackslash}p{6.8cm}|>{\raggedright\arraybackslash}p{6.8cm}|}
		\hline
		\rowcolor{eccblue!15}
		\textbf{\color{eccblue}\#} & \textbf{\color{eccblue}Écart constaté} &
		\textbf{\color{eccblue}Correction apportée} \\
		\hline
		\rowcolor{white}
		1 & TRI et VAN affichés en tête & Rendement net et cash-flow remontés en
		première position sur tous les supports \\
		\hline
		\rowcolor{gray!8}
		2 & Un loyer supposé dans tous les cas & Opération sans loyer traitée
		nativement, valeur créée par la plus-value \\
		\hline
		\rowcolor{white}
		3 & Facteur temps secondaire & Horizon et délai de livraison devenus des
		paramètres de premier plan \\
		\hline
		\rowcolor{gray!8}
		4 & Phase amont absente & Brief de recherche ajouté au modèle de données \\
		\hline
		\rowcolor{white}
		5 & Fiche client incomplète & Situation professionnelle, nationalité et
		budget disponible ajoutés \\
		\hline
		\rowcolor{gray!8}
		6 & Déroulé du dossier non suivi & Six statuts d'avancement ajoutés \\
		\hline
		\rowcolor{white}
		7 & VEFA traitée comme un bien louable & Délai de livraison modélisé et
		proratisé \\
		\hline
	\end{tabular}
	\caption{Les sept écarts relevés entre le cahier des charges initial et la
		pratique du cabinet.}
\end{table}

\subsection{Points laissés hors périmètre}

Signalés comme utiles par le cabinet mais non réalisés, faute de temps avant la
remise. Ils sont documentés pour que la reprise du projet parte d'un état connu :

\begin{itemize}
	\item \textbf{Échéancier VEFA par tranches} : les appels de fonds suivent
	l'avancement du chantier ; l'outil modélise le décalage de livraison, pas
	l'étalement des paiements.
	\item \textbf{Calcul de capacité d'emprunt} à partir de la situation
	professionnelle et du budget disponible.
	\item \textbf{Rapprochement automatique} entre le brief d'un client et un
	portefeuille de biens : l'outil enregistre le besoin, il ne cherche pas encore
	les biens.
	\item \textbf{Suivi des visites} (dates, retours du client) au sein de l'étape
	correspondante.
	\item \textbf{Validation formelle des taux par défaut} des zones de marché
	(Maroc : frais 7 \%, imposition effective 15 \%), aujourd'hui des valeurs de
	travail documentées comme telles.
\end{itemize}

\subsection{Liste des informations à compléter avant envoi}

\begin{encadre}
	Les éléments suivants apparaissent en rouge dans le corps du rapport et doivent
	être renseignés avant l'envoi :
	\begin{itemize}\itemsep0.15em
		\item numéro d'étudiant, en tête du fichier source ;
		\item nom, prénom et fonction du tuteur, en tête du fichier source ;
		\item modalité du stage (présentiel, hybride, distanciel) et rythme de
		présence ;
		\item date de création, forme juridique et effectif du cabinet ;
		\item service ou entité de rattachement et nombre de personnes y travaillant ;
		\item position concurrentielle du cabinet sur le Grand Casablanca ;
		\item noms des collaborateurs rencontrés, pour les remerciements ;
		\item intitulés exacts des cours de première année mobilisés.
	\end{itemize}
	Le marqueur \verb|\acompleter{...}| du fichier source signale chacun d'eux ;
	une recherche sur ce mot-clé permet de tous les retrouver.
\end{encadre}

\end{document}
